% TOP_PROBLEM: {"id": 2305073, "title": "Research Problems in Function Theory — Problem 5.73", "category_id": 8, "category": {"id": 8, "name": "analysis", "display_name": "Analysis", "description": "Limits, continuity, calculus, and function theory.", "slug": "analysis", "order_index": 8, "created_at": "2026-07-31T15:26:25.671Z"}, "status": "open", "problem_number": "AMR-022-5073", "set_id": 13}
% FIRST_POSTED: 2026-10-02
% ENTRY_KIND: statement_only
% UnsolvedMath ID 2305073; source catalogue code AMR-022-5073
% !TeX program = lualatex
% Definition and sources only; this file is not an LLM solution attempt.
\documentclass[11pt,a4paper]{article}
\usepackage[margin=25mm]{geometry}
\usepackage{fontspec}
\setmainfont{lmroman10-regular.otf}[BoldFont=lmroman10-bold.otf,
 ItalicFont=lmroman10-italic.otf,BoldItalicFont=lmroman10-bolditalic.otf]
\usepackage{amsmath,amssymb,mathtools}
\usepackage{unicode-math}
\setmathfont{latinmodern-math.otf}
\AtBeginDocument{\let\setminus\smallsetminus}
\usepackage{xurl,hyperref}
\hypersetup{colorlinks=true,urlcolor=blue}
\setcounter{secnumdepth}{0}
\setlength{\parindent}{0pt}
\setlength{\parskip}{5pt}
\setlength{\emergencystretch}{3em}
\begin{document}
\section*{Research Problems in Function Theory — Problem 5.73}\label{2305073}
{\small UnsolvedMath ID 2305073 · AMR-022-5073 · Analysis · AMR Open Problem Lists}
\subsection{Definitions and mathematical statement}
Fix $0<\alpha<1$. For a finite positive Borel measure $\mu$ on $\mathbb R$, define its Riesz potential by
\[
I_\alpha\mu(x)=\int_{\mathbb R}|x-t|^{-\alpha}\,d\mu(t),
\]
allowing the value $+\infty$. Here ``finite'' means $\mu(\mathbb R)<\infty$; neither compact support nor absolute continuity is required. In particular atoms are allowed, and the potential diverges at any atom.
Characterize the nonnegative measurable functions $g$ on $\mathbb R$ for which there exists such a measure satisfying
\[
g(x)\le I_\alpha\mu(x)\qquad\text{for every }x\in\mathbb R.
\]
This pointwise formulation follows the stated domination question; domination only almost everywhere is a distinct possible variant and should be identified if used. A useful criterion should be expressed directly in terms of $g$, for instance through integral, covering, or capacity conditions, and should be both necessary and sufficient.
The total mass of the witnessing measure may depend on $g$, and the normalization of the kernel can be absorbed into that mass. The problem does not require $g$ itself to be a potential, only to lie below one. This difference permits highly irregular measurable functions while preserving the global finite-mass constraint on their majorants.
\subsection{Short English statement}
Which nonnegative measurable functions on the line admit a pointwise majorant that is a Riesz potential of a finite positive measure?
\subsection{Sources}
\begin{itemize}
\item[S1] ULAM.AI and the UnsolvedMath Contributors, supplied problems.json, record 2305073 (AMR-022-5073), AMR Open Problem Lists. Catalogue identity and source transcription; CC BY 4.0.
\url{https://huggingface.co/datasets/ulamai/UnsolvedMath}.
\item[S2] Hayman - Research Problems in Function Theory (2018), Problem 5.73. Original source indexed by AMR-022-5073.
\url{https://arxiv.org/abs/1809.07200}.
\end{itemize}
\end{document}
